\section{8. Beispiele \texorpdfstring{$\star$}{*}}

\subsection{Grenzwerte}
\Bsp \key{Trick mit Wurzeln}: Lösung von $\lim\limits_{n\to\infty} \sqrt{n^2 + 3n} - n$
\begin{align*}
    \lim\limits_{n\to\infty} \sqrt{n^2 + 3n} - n &=\lim\limits_{n\to\infty} (\sqrt{n^2 + 3n} - n) \cdot \frac{\sqrt{n^2 + 3n} + n}{\sqrt{n^2 + 3n} + n} \\[-0.6em]
        &= \lim\limits_{n\to\infty} \frac{n^2 + 3n - n^2}{\sqrt{n^2 + 3n} + n} \\[-0.2em]
        &= \lim\limits_{n\to\infty} \frac{3n}{\sqrt{n^2 + 3n} + n} \\[-.3em]
        &= \lim\limits_{n\to\infty} \frac{3}{\sqrt{1 + \frac{3}{n}} + 1} \\[-.3em]
        &= \frac{3}{\sqrt{1 + 0} + 1} = \frac{3}{2}
\end{align*}

\subsection{Differentialgleichungen}

\Bsp \key{Trennung der Variablen}: Lösung von $y' = \frac{e^x(y^2 - 1)}{y(e^x + 1)}$
\vspace{-0.2em}
\begin{align*}
    y' &= \frac{e^x(y^2 - 1)}{y(e^x + 1)} \\[-0.4em]
       &= \frac{e^x}{e^x + 1} \cdot \frac{y^2 - 1}{y} \\[-0.4em]
    \overset{TdV}{\implies} \int \frac{y}{y^2 - 1}dy &= \int \frac{e^x}{e^x + 1}dx \\[-0.3em]
    \int \frac{1}{2} \cdot \frac{1}{u} du &= \int \frac{1}{u} du \\[-0.5em]
    \frac{1}{2}\ln(|y^2 - 1|) &= \ln(|e^x + 1|) + C_1 \\[-0.6em]
    \ln(|y^2 - 1|) &= 2\ln(|e^x + 1|) + C_2 \\[-0.6em]
    |y^2 - 1| &= (e^x + 1)^2 \cdot C_3 \\[-0.6em]
    y^2 - 1 &= \pm C_3 \cdot (e^x + 1)^2 \\[-0.6em]
            &= C_4 \cdot (e^x + 1)^2 \\[-0.6em]
    y^2 &= C_4 (e^x + 1)^2 + 1 \\[-0.6em]
    y &= \pm \sqrt{C_4 (e^x + 1)^2 + 1}
\end{align*}

\begin{minipage}{\linewidth}
\Bsp \key{Substitution}: Lösung von $y' = y + x$
\begin{enumerate}
    \item Substitution $u = x + y$, also $u' = y' + 1$ (d.h. $y' = u' - 1$)
    \item Setze in der DGL ein: $u' - 1 = u \implies u' = u + 1$
    \item Löse durch Trennung der Variablen: $u = Ce^x - 1$
    \item Setze $u = x + y$ ein: $x + y = Ce^x - 1 \implies y = Ce^x - x - 1$
\end{enumerate}
\end{minipage}

\begin{minipage}{\linewidth}
\Bsp \key{Variation der Konstanten}: Lösung von $y' - 2y = e^{3x}$
\begin{enumerate}
    \item Löse $y_h' - 2y_h = 0$:
        \begin{align*}
            y_h' - 2y_h &= 0 \\[-0.6em]
            y_h' &= 2y_h \\[-0.6em]
            y_h &= Ce^{2x} \\[-0.6em]
        \end{align*}
    \item Betrachte Konstanten als Funktion: $C \to C(x)$, $y = C(x) \cdot e^{2x}$
    \item Berechne Ableitung(en): $y' = C'(x) \cdot e^{2x} + 2 C(x) \cdot e^{2x}$
    \item Berechne $C(x)$:
        \vspace{-0.6em}
        \begin{align*}
            y' - 2y &= e^{3x} \\[-0.6em]
            C'(x) \cdot e^{2x} + 2 C(x) \cdot e^{2x} - 2C(x) \cdot e^{2x} &= e^{3x} \\[-0.6em]
            C'(x) \cdot e^{2x} &= e^{3x} \\[-0.6em]
            C'(x) &= e^x \\[-0.6em]
            C(x) &= e^x + C \\[-0.6em]
        \end{align*}
    \item Berechne $y = C(x) \cdot e^{2x} = (e^x + C) \cdot e^{2x} = e^{3x} + Ce^{2x}$
\end{enumerate}
\end{minipage}

\Bsp \key{Ansätze für die partikuläre Lösung}: Lösung von $y''' - 2y'' = 4e^{2x} + 12$
\begin{enumerate}
    \item Löse $y_h''' - 2y_h'' = 0$ (char. Polynom $\lambda^3 - 2\lambda^2 = \lambda^2(\lambda - 2)$)
        \vspace{-0.4em}
        \begin{align*}
            y_h &= \sum_{i = 1}^{2} (\sum_{p=0}^{m_i - 1} C_{ip} x^p e^{\lambda_i x}) \\[-0.6em]
                &= \sum_{p = 0}^1 C_{1p} x^p + C_{20} e^{2x} \\[-0.6em]
                &= C_{10} + C_{11}x + C_{20}e^{2x} \\[-0.6em]
                &= C_1 + C_2x + C_3e^{2x} \\[-0.6em]
        \end{align*}
    \item Finde $y_p$ und Ableitungen mit Ansätze (beide sind Spezialfälle):
        \vspace{-1em}
        \begin{align*}
            y_p &= C_1 x e^{2x} + C_2 x^2 \\[-0.4em]
            y_p' &= C_1 e^{2x} + 2 C_1 x e^{2x} + 2C_2 x \\[-0.4em]
            y_p'' &= 4C_1e^{2x} + 4C_1 x e^{2x} + 2C_2 \\[-0.4em]
            y_p''' &= 12 C_1 e^{2x} + 8C_1 x e^{2x}
        \end{align*}
    \item Substitution in der DGL um $C_1, C_2$ zu bestimmen:
        \begin{align*}
            y_p''' - 2y_p'' &= 4e^{2x} + 12 \\[-0.3em]
            12C_1 e^{2x} + 8C_1xe^{2x} - 8C_1e^{2x} - 8C_1xe^{2x} - 4C_2 &= 4e^{2x} + 12 \\[-0.6em]
            4C_1 e^{2x} - 4C_2 &= 4e^{2x} + 12 \\[-0.6em]
            \implies C_1 = 1, C_2 = -3 \hspace{60pt}& \\[-0.6em]
            \implies y_p = x e^{2x} - 3 x^2 \hspace{61pt}&
        \end{align*}
    \item Allgemeine Lösung $y = y_h + y_p = C_1 + C_2x + C_3e^{2x} + x e^{2x} - 3 x^2$
\end{enumerate}

\Bsp \key{Anfangswertproblem}: Lösung von $y'' - 6y' + 25y = 0$ gegeben $y(0) = 0$ und $y'(0) = 1$
\begin{enumerate}
    \item Löse DGL (char. Polynom $\lambda^2 - 6\lambda + 25$ hat Lösungen $\lambda_{1, 2} = 3 \pm 4i$)
    \begin{align*}
        y &= \sum_{j=1}^{s} \left(\sum_{q=0}^{m_j-1}(A_{jq} x^q e^{a_j x} \sin(b_j x) + B_{jq} x^q e^{a_j x} \cos(b_j x))\right) \\
          &= e^{3x} (A \sin(4 x) + B  \cos(4 x))
    \end{align*}
    \item Berechne $y'(x)$: $y'(x) = e^{3x} ((3A - 4B) \sin(4x) + (4A + 3B) \cos(4x))$
    \item Berechne $y(0) = B$ und $y'(0) = 4A + 3B$
    \item Setze $y(0) = 0$ und $y'(0) = 1$ und Löse das System: $B = 0$, $A = \frac{1}{4}$,
     also  $y(x) = \frac{1}{4}e^{3x} \sin(4 x)$
\end{enumerate}

\subsection{Integration}

\Bsp \key{Partielle Integration}: Lösung von $\displaystyle \int e^{2x} \cdot \cos(x) dx$
\vspace{-0.6em}
\begin{align*}
     I &= \int \underbrace{e^{2x}}_\downarrow \underbrace{\cos x}_\uparrow dx \\[-0.8em]
       &= e^{2x} \sin x - 2 \int \underbrace{e^{2x}}_\downarrow \underbrace{\sin x}_\uparrow dx \\[-0.8em]
       &= e^{2x} \sin x - 2 (e^{2x} \cos x + 2 \int e^{2x} \cos x dx) \\[-0.6em]
       &= e^{2x} \sin x + 2 e^{2x} \cos x - 4 I \\[-0.2em]
    \implies 5I &= e^{2x} \sin x + 2 e^{2x} \cos x \\[-0.4em]
    \implies \int e^{2x} \cdot \cos(x) dx &= \frac{1}{5}e^{2x}(\sin x + 2\cos x) + C
\end{align*}


\begin{minipage}{\linewidth}
\Bsp \key{DI-Methode} (mehrfache part. Int. als Tabelle): Lösung von $\displaystyle \int x^2 e^x dx$
\vspace{-0.6em}
\begin{center}
\begin{tabular}{c|c c}
     & $D$ (ableiten) & $I$ (integrieren) \\
    \hline
    $+$ & $x^2$ & $e^x$ \\
    $-$ & $2x$ & $e^x$ \\
    $+$ & $2$ & $e^x$ \\
     & $0$ &
\end{tabular}
\end{center}
\vspace{-0.4em}
Diagonalprodukte mit Vorzeichen aufsummieren (letzte Zeile mit $D = 0$ weglassen):
$$\int x^2 e^x dx = x^2e^x - 2xe^x + 2e^x + C$$
\end{minipage}

\begin{minipage}{\linewidth}
\Bsp \key{Partielle Integration + Subtitution}: Lösung von $\displaystyle \int \arcsin(x) dx$
\vspace{-0.6em}
\begin{align*}
    \int \arcsin(x) dx &= \int \underbrace{1}_\uparrow \cdot \underbrace{\arcsin(x)}_\downarrow dx \\[-2.6em]
                       &= x \cdot \arcsin(x) - \int x \cdot \overbrace{\frac{1}{\sqrt{1-x^2}}}^{\frac{1}{\sqrt{y}}} dx \\
    \longrightarrow y = 1 - x^2 \implies& \frac{dy}{dx} = -2x \implies \frac{-1}{2} \, dy = x \, dx \\[-0.1em]
    \implies \int x \cdot \frac{1}{\sqrt{1-x^2}} dx &= \int \frac{-1}{2} \cdot \frac{1}{\sqrt{y}} dy \\[-0.1em]
                       &= -\frac{1}{2} \cdot \int y^{-\frac{1}{2}} dy \\[-0.3em]
                       &= -\frac{1}{2} (2y^{\frac{1}{2}} + C) \\[-0.3em]
                       &= -\sqrt{y} + C = -\sqrt{1-x^2} + C \\
    \implies \int \arcsin(x) dx &= x \cdot \arcsin(x) - (-\sqrt{1 - x^2} + C) \\[-0.8em]
                                &= x \cdot \arcsin(x) + \sqrt{1 - x^2} + C
\end{align*}
\end{minipage}

\Bsp \key{Partialbruchzerlegung}: Lösung von $\displaystyle \int \frac{x^2 + 2x - 1}{x^3 - x} dx$
\vspace{-0.6em}
\begin{enumerate}
    \item $2 < 3$, also brauchen wir keine division
    \item $x^3 - x = x(x+1)(x-1)$
    \item $g(x) = \frac{A}{x} + \frac{B}{x + 1} + \frac{C}{x - 1}$
    \item $\frac{x^2 + 2x - 1}{x^3 - x} = \frac{A}{x} + \frac{B}{x + 1} + \frac{C}{x - 1} \\
    \implies x^2 + 2x -1 = A(x^2 - 1) + B(x^2 - x) + C(x^2 + x)$
    \item Lösen durch einsetzen $x = 0, \pm 1$: $A = C = 1, B = -1$
    \item $\displaystyle \int \frac{1}{x} + \frac{-1}{x+1} + \frac{1}{x-1} dx = \ln(|x|) - \ln(|x+1|) + \ln(|x-1|) + C$
\end{enumerate}

\Bsp \key{Weierstrass-Substitution}: Lösung von $\displaystyle \int \frac{1}{2 + \sin x} dx$:
\begin{align*}
    t = \tan(\frac{x}{2}), \hspace{10pt} \sin x &= \frac{2t}{1 + t^2}, \hspace{10pt} dx = \frac{2}{1 + t^2} dt \\[-0.1em]
    \implies \int \frac{1}{2 + \sin x} dx &= \int \frac{1}{2 + \frac{2t}{1 + t^2}} \cdot \frac{2}{1 + t^2} dt \\[-0.1em]
        &= \int \frac{2}{\frac{2 + 2t^2 + 2t}{1+t^2} \cdot (1+t^2)} dt \\[-0.1em]
        &= \int \frac{1}{(t + \frac{1}{2})^2 + \frac{3}{4}} dt \\[-0.4em]
        &= \frac{2}{\sqrt{3}} \arctan(\frac{t + \frac{1}{2}}{\frac{\sqrt{3}}{2}}) + C \\[-0.4em]
        &= \frac{2}{\sqrt{3}} \arctan(\frac{\tan(\frac{x}{2}) + \frac{1}{2}}{\frac{\sqrt{3}}{2}}) + C
\end{align*}

\subsection{Funktionenfolgen}

\begin{minipage}{\linewidth}
\Bsp \key{Punktw. vs. gleichm. Konvergenz}: Untersuche $f_n(x) = x^n$ auf $[0,1]$
\begin{enumerate}
    \item \key{Punktweise Konvergenz}: Für $x \in [0,1)$: $\lim\limits_{n\to\infty} x^n = 0$.
    Für $x=1$: $\lim\limits_{n\to\infty} 1^n = 1$
    \item \key{Grenzfunktion}: $f(x) = \begin{cases} 0 & x \in [0,1) \\ 1 & x = 1 \end{cases}$
    \item \key{Gleichmässige Konvergenz}: $M_n = \sup\limits_{x\in[0,1]}|x^n - f(x)| = \sup\limits_{x\in[0,1)} x^n = 1$
    (da $x^n \to 1$ für $x \to 1^-$), also $\limsup\limits_{n\to\infty} M_n = 1 \ne 0$
    $\implies$ \key{nicht} gleichmässig konvergent
\end{enumerate}
Beachte: $f_n$ stetig, aber $f$ unstetig bei $x=1$ -- passt zum Satz, dass gleichm.
Konv. stetiger Funkt. wieder stetig ist
\end{minipage}

\subsection{Misc}

\Bsp \key{Polynomdivision}: $p(x) = s(x) \cdot q(x) + r(x)$
\begin{align*}
    &\overbrace{7x^4 \hspace{7pt} + 3x^2 - 6x + 1}^{p(x)} : \overbrace{4x^2 + 3}^{q(x)} = \overbrace{\frac{7}{4}x^2 - \frac{9}{16}}^{s(x)} \\[-0.6em]
    &\hspace{-9pt}-7x^4 - \frac{21}{4}x^2 \\
    &\hspace{19pt} -\frac{9}{4}x^2 - 6x + 1 \\[-0.3em]
    &\hspace{19pt} +\frac{9}{4}x^2 \hspace{19.5pt} + \frac{27}{16} \\[-0.2em]
    &\hspace{44pt} \left. -6x \hspace{2.5pt} + \frac{43}{16} \right\} r(x)
\end{align*}
